\documentclass[12pt]{ximera}
\title{Practice Exam 1}
\begin{document}
\begin{abstract}
Practice for Exam 1, covering Chapter 1: Condorcet candidates and the Pairwise
Comparisons method, a preference schedule with no Condorcet candidate, bookkeeping for
Borda points, and Plurality-with-Elimination.
\end{abstract}
\maketitle

\section*{Practice Exam 1}

\begin{problem}
Assess whether each of the following claims is true or false.

\begin{prompt}
\textbf{(a)} Claim: it is possible for a Condorcet candidate to lose an election using
the Pairwise Comparisons method.

\begin{multipleChoice}
\choice{True}
\choice[correct]{False}
\end{multipleChoice}

\begin{feedback}
False. A Condorcet candidate wins every head-to-head comparison, so earns the maximum
possible points --- more than anyone else can.
\end{feedback}
\end{prompt}

\begin{prompt}
\textbf{(b)} Claim: it is possible to win an election using the Pairwise Comparisons
method without being a Condorcet candidate.

\begin{multipleChoice}
\choice[correct]{True}
\choice{False}
\end{multipleChoice}

\begin{feedback}
True. The winner only needs the most points, not a clean sweep. When no Condorcet
candidate exists --- as in Problem 2 --- someone still wins, despite losing at least one
matchup.
\end{feedback}
\end{prompt}
\end{problem}

\begin{problem}
Consider the following preference schedule for an election with four candidates.

\begin{center}
\begin{tabular}{c|ccc}
Number of votes & 8 & 6 & 4 \\\hline
1st & B & A & D \\
2nd & C & D & B \\
3rd & A & C & C \\
4th & D & B & A
\end{tabular}
\end{center}

\begin{prompt}
\textbf{(a)} Give the complete ranking by the Plurality method.

1st: $\answer{B}$ with $\answer{8}$ \quad 2nd: $\answer{A}$ with $\answer{6}$ \quad
3rd: $\answer{D}$ with $\answer{4}$ \quad 4th: $\answer{C}$ with $\answer{0}$

\begin{feedback}
First-place votes: B $8$, A $6$, D $4$, C $0$.
\end{feedback}
\end{prompt}

\begin{prompt}
\textbf{(b)} Is there a Condorcet candidate in this election?

\begin{multipleChoice}
\choice{Yes, A}
\choice{Yes, B}
\choice{Yes, C}
\choice{Yes, D}
\choice[correct]{No}
\end{multipleChoice}

\begin{feedback}
No. The head-to-head results are
\[
\begin{array}{lll}
\text{A vs B} & 6\text{--}12 & \text{B}\\
\text{A vs C} & 6\text{--}12 & \text{C}\\
\text{A vs D} & 14\text{--}4 & \text{A}\\
\text{B vs C} & 12\text{--}6 & \text{B}\\
\text{B vs D} & 8\text{--}10 & \text{D}\\
\text{C vs D} & 8\text{--}10 & \text{D}
\end{array}
\]
B loses to D, D loses to A, and A loses to B --- everyone loses at least once, so no one
beats every rival.
\end{feedback}
\end{prompt}
\end{problem}

\begin{problem}
Suppose an election with five candidates (A, B, C, D, E) and $100$ voters is run using
the standard Borda Count method.

\begin{prompt}
\textbf{(a)} How many Borda points is each ballot worth?

$\answer{15}$

\begin{feedback}
A ballot awards $5+4+3+2+1=15$ points.
\end{feedback}
\end{prompt}

\begin{prompt}
\textbf{(b)} How many total Borda points are awarded?

$\answer{1500}$

\begin{feedback}
$100$ ballots at $15$ points each: $1500$.
\end{feedback}
\end{prompt}

\begin{prompt}
\textbf{(c)} Candidate A earns $300$ Borda points, B earns $250$, C earns $350$, and D
earns $200$. How many did E earn?

$\answer{400}$

\begin{feedback}
$1500-(300+250+350+200)=1500-1100=400$.
\end{feedback}
\end{prompt}

\begin{prompt}
\textbf{(d)} Give the complete ranking by Borda count.

1st $\answer{E}$ \quad 2nd $\answer{C}$ \quad 3rd $\answer{A}$ \quad 4th $\answer{B}$
\quad 5th $\answer{D}$

\begin{feedback}
E $400$, C $350$, A $300$, B $250$, D $200$.
\end{feedback}
\end{prompt}
\end{problem}

\begin{problem}
Consider the following preference schedule for an election with five candidates.

\begin{center}
\begin{tabular}{c|ccccc}
Number of votes & 10 & 8 & 6 & 4 & 2\\\hline
1st & D & B & A & E & C \\
2nd & B & A & B & D & D \\
3rd & E & C & D & B & A \\
4th & A & D & E & C & E \\
5th & C & E & C & A & B
\end{tabular}
\end{center}

\begin{prompt}
\textbf{(a)} How many voters participated? $\answer{30}$

\begin{feedback}
$10+8+6+4+2=30$, so a majority is $16$ votes.
\end{feedback}
\end{prompt}

\begin{prompt}
\textbf{(b)} Using Plurality-with-Elimination, which candidate is eliminated in Round 1?

$\answer{C}$

\begin{feedback}
First-place votes: D $10$, B $8$, A $6$, E $4$, C $2$. C has the fewest.
\end{feedback}
\end{prompt}

\begin{prompt}
\textbf{(c)} Which candidate is eliminated in Round 2?

$\answer{E}$

\begin{feedback}
C's $2$ voters ranked D next, so D rises to $12$: D $12$, B $8$, A $6$, E $4$. E is
eliminated.
\end{feedback}
\end{prompt}

\begin{prompt}
\textbf{(d)} Does the election require further elimination to determine the winner?

\begin{multipleChoice}
\choice{Yes --- no candidate has a majority yet}
\choice[correct]{No --- D now has a majority}
\end{multipleChoice}

\begin{feedback}
No. E's $4$ voters ranked D next, so D rises to $16$ of the $30$ votes --- a majority.
D is the winner, and further rounds are needed only to rank the others.
\end{feedback}
\end{prompt}

\begin{prompt}
\textbf{(e)} Give the complete ranking.

1st $\answer{D}$ \quad 2nd $\answer{B}$ \quad 3rd $\answer{A}$ \quad 4th $\answer{E}$
\quad 5th $\answer{C}$

\begin{feedback}
Continuing to rank the rest: A ($6$) is eliminated next, and A's voters go to B, giving
D $16$, B $14$. Reversing the elimination order: D, B, A, E, C.
\end{feedback}
\end{prompt}
\end{problem}

\end{document}
